\subsection{Sustitución: construcción simple de un k-d tree}
\label{sec:act3-2}

\begin{enunciado}
Considere una construcción de k-d tree en la cual, en cada llamada recursiva, los puntos
se ordenan por la coordenada correspondiente al nivel, se selecciona la mediana y se
construyen recursivamente ambos subárboles.
\begin{lstlisting}[language={}, numbers=none, frame=none, backgroundcolor={}]
BuildKD(P, depth)
    if |P| <= 1
        return leaf(P)

    axis = depth mod d
    sort P according to coordinate axis
    m = floor(|P|/2)

    node.point = P[m]
    node.left  = BuildKD(P[0 ... m-1], depth+1)
    node.right = BuildKD(P[m+1 ... |P|-1], depth+1)
    return node
\end{lstlisting}
Suponga que ordenar $m$ puntos tiene costo $\T{m \log m}$ y que las particiones quedan
balanceadas.
\begin{enumerate}[label=(\alph*)]
  \item Obtenga la recurrencia correspondiente.
  \item Proponga una cota asintótica para la recurrencia. Puede utilizar un árbol de
        recursión únicamente para formular la conjetura.
  \item Demuestre la cota mediante el método de sustitución. La demostración debe incluir
        cota superior y cota inferior, o una argumentación equivalente que permita concluir
        una cota $\Theta$.
  \item Si en vez de ordenar completamente en cada llamada la mediana pudiera obtenerse y
        particionarse en $\T{m}$, escriba la nueva recurrencia y determine cómo cambiaría
        el costo de construcción.
\end{enumerate}
\end{enunciado}

% ============================================================
\subsubsection*{(a) Recurrencia}

Llamamos $T(n)$ al tiempo que tarda \texttt{BuildKD} con un conjunto $P$ de $n$ puntos.
Recorremos el pseudocódigo línea por línea.

Si $n \le 1$, el algoritmo compara $|P|$ con $1$ y crea una hoja. Eso cuesta una
constante, que llamamos $c$.

Si $n \ge 2$, el algoritmo compara $|P|$ con $1$, calcula \texttt{axis} y \texttt{m}, y
guarda \texttt{node.point}. Cada una de estas operaciones cuesta $\T{1}$. Además ordena
los $n$ puntos según la coordenada \texttt{axis}, lo que según el enunciado cuesta
$\T{n \log n}$. Juntando todo, el trabajo que no es recursivo es $\T{n \log n}$, porque
las operaciones de costo constante quedan absorbidas por el ordenamiento. Al final hay
dos llamadas recursivas: una con los puntos $P[0 \dots m-1]$ y otra con
$P[m+1 \dots |P|-1]$. Como $m = \floor{n/2}$, cada una recibe a lo más $n/2$ puntos, y
como las particiones quedan balanceadas, cada llamada tarda $T(n/2)$. Suponemos que
pasar un subarreglo no copia los puntos, así que cuesta $\T{1}$.

Para no cargar la notación, usamos la misma constante $c$ para el caso base y para el
trabajo no recursivo, que escribimos como $c\,n\log_2 n$. Entonces:
\[
  T(n) =
  \begin{cases}
    c, & n \le 1, \\[4pt]
    2\,T\!\left(\dfrac{n}{2}\right) + c\,n \log_2 n, & n \ge 2.
  \end{cases}
\]
El $2\,T(n/2)$ viene de las dos llamadas recursivas, \texttt{node.left} y
\texttt{node.right}, y el término $c\,n\log_2 n$ viene del ordenamiento. Como se permite
en clases, suponemos que $n$ es una potencia de $2$.

% ============================================================
\subsubsection*{(b) Conjetura con el árbol de recursión}

Usamos el árbol de recursión para formular la conjetura, que luego se demuestra por
sustitución en (c). Cada nodo es una llamada. Dentro de cada
nodo anotamos el tamaño del subproblema y, debajo, su trabajo no recursivo. La raíz tiene
tamaño $n$ y trabajo $c\,n\log_2 n$, y cada nodo de tamaño $s$ tiene dos hijos de tamaño
$s/2$.

\begin{center}
\begin{forest}
  for tree={
    draw, rounded corners, align=center, font=\small,
    l sep=12mm, s sep=3mm, inner sep=3pt, edge={->}
  }
  [{$n$\\[1pt]$c\,n\log_2 n$}
    [{$n/2$\\[1pt]$c\,\frac{n}{2}\log_2\frac{n}{2}$}
      [{$n/4$\\[1pt]$c\,\frac{n}{4}\log_2\frac{n}{4}$}]
      [{$n/4$\\[1pt]$c\,\frac{n}{4}\log_2\frac{n}{4}$}]
    ]
    [{$n/2$\\[1pt]$c\,\frac{n}{2}\log_2\frac{n}{2}$}
      [{$n/4$\\[1pt]$c\,\frac{n}{4}\log_2\frac{n}{4}$}]
      [{$n/4$\\[1pt]$c\,\frac{n}{4}\log_2\frac{n}{4}$}]
    ]
  ]
\end{forest}
\end{center}

En el nivel $i$ hay $2^i$ nodos, cada uno de tamaño $n/2^i$. Los tamaños se dividen por
$2$ en cada nivel, así que llegan a $1$ en el nivel $\log_2 n$, donde están las $n$
hojas. Cada hoja cuesta $c$.
\begin{center}
\begin{tabular}{cccc}
  \toprule
  nivel & nodos & tamaño de cada nodo & trabajo del nivel \\
  \midrule
  0 & $1$   & $n$     & $c\,n \log_2 n$ \\[2pt]
  1 & $2$   & $n/2$   & $2 \cdot c\,\frac{n}{2} \log_2 \frac{n}{2} = c\,n\,(\log_2 n - 1)$ \\[4pt]
  2 & $4$   & $n/4$   & $4 \cdot c\,\frac{n}{4} \log_2 \frac{n}{4} = c\,n\,(\log_2 n - 2)$ \\[2pt]
  $\vdots$ & $\vdots$ & $\vdots$ & $\vdots$ \\
  $i$ & $2^i$ & $n/2^i$ & $c\,n\,(\log_2 n - i)$ \\
  $\vdots$ & $\vdots$ & $\vdots$ & $\vdots$ \\
  $\log_2 n - 1$ & $n/2$ & $2$ & $c\,n \cdot 1$ \\[2pt]
  $\log_2 n$ (hojas) & $n$ & $1$ & $c\,n$ \\
  \bottomrule
\end{tabular}
\end{center}
Sumando los niveles internos, del $0$ al $\log_2 n - 1$, el factor entre paréntesis toma
los valores $\log_2 n, \log_2 n - 1, \dots, 2, 1$:
\[
  c\,n\,\big(\log_2 n + (\log_2 n - 1) + \cdots + 2 + 1\big)
  = c\,n \cdot \frac{\log_2 n\,(\log_2 n + 1)}{2}
  = \frac{c}{2}\, n \log_2^2 n + \frac{c}{2}\, n \log_2 n,
\]
donde usamos la suma $1 + 2 + \cdots + L = L(L+1)/2$ con $L = \log_2 n$. Agregando las
hojas, el árbol suma
\[
  \frac{c}{2}\, n \log_2^2 n + \frac{c}{2}\, n \log_2 n + c\,n.
\]
El término que más crece es $\frac{c}{2}\, n\log_2^2 n$. A diferencia de la Actividad
3.1, el trabajo de los niveles no baja en forma geométrica: baja solo de a $c\,n$ por
nivel, así que la mitad de los niveles cuesta al menos la mitad de lo que cuesta la raíz.
Son $\log_2 n$ niveles que cuestan del orden de $n\log_2 n$ cada uno, y por eso
conjeturamos
\[
  T(n) \in \T{n \log^2 n}.
\]
Es el mismo ejemplo que se vio en clases para la forma general del caso 2 del Teorema
Maestro: $2T(n/2) + n\log n \in \T{n\log^2 n}$.

% ============================================================
\subsubsection*{(c) Demostración por sustitución}

En las dos cotas usamos que, como $n$ es potencia de $2$,
\[
  \log_2 \frac{n}{2} = \log_2 n - 1,
  \qquad \text{y por lo tanto} \qquad
  \left( \log_2 \frac{n}{2} \right)^2 = \log_2^2 n - 2\log_2 n + 1.
\]

\textbf{Cota superior.} Conjetura: $T(n) \le 2c\, n \log_2^2 n$ para todo $n \ge 2$.

\emph{Caso base, $n = 2$.} Por la recurrencia, $T(2) = 2\,T(1) + c \cdot 2 \cdot \log_2 2$.
Como $T(1) = c$ y $\log_2 2 = 1$, queda $T(2) = 2c + 2c = 4c$. La cota en $n = 2$ vale
$2c \cdot 2 \cdot (\log_2 2)^2 = 4c$, así que se cumple $4c \le 4c$. Este caso también
explica la constante. Si la conjetura fuera $T(n) \le d\, n\log_2^2 n$, en $n = 2$ se
necesitaría $4c \le 2d$, es decir, $d \ge 2c$, y tomamos el menor valor posible,
$d = 2c$.

La base no puede ser $n = 1$, porque ahí la cota vale $2c \cdot 1 \cdot (\log_2 1)^2 = 0$
y $T(1) = c$ es mayor que $0$. La cota es falsa en $n = 1$, y por eso la conjetura se
plantea desde $n = 2$. Partir de $n = 2$ alcanza para cubrir todos los casos. Como $n$ es
potencia de $2$, los valores que importan son $2, 4, 8, 16, \dots$, y cada uno es el doble
del anterior. El paso inductivo lleva la cota de $n/2$ a $n$, así que desde $n = 2$ se
llega a $4$, de $4$ a $8$, y así sucesivamente. El paso parte en $n = 4$ para que $n/2$
sea al menos $2$ y la hipótesis se pueda usar.

\emph{Paso inductivo.} Sea $n \ge 4$ y supongamos que la cota vale para $n/2$, es decir,
$T(n/2) \le 2c\, \frac{n}{2} \log_2^2 \frac{n}{2}$. Sustituyendo en la recurrencia:
\[
\begin{aligned}
  T(n) &= 2\,T\!\left(\frac{n}{2}\right) + c\,n\log_2 n \\
       &\le 2 \cdot 2c\, \frac{n}{2} \left(\log_2 n - 1\right)^2 + c\,n\log_2 n \\
       &= 2c\,n \left(\log_2^2 n - 2\log_2 n + 1\right) + c\,n\log_2 n \\
       &= 2c\,n\log_2^2 n - 3c\,n\log_2 n + 2c\,n.
\end{aligned}
\]
Para que esto sea a lo más $2c\,n\log_2^2 n$, lo que sobra tiene que ser negativo o cero:
$-3c\,n\log_2 n + 2c\,n \le 0$. Dividiendo por $c\,n$, eso es $3\log_2 n \ge 2$, que se
cumple porque $\log_2 n \ge 2$ para $n \ge 4$. Por lo tanto, $T(n) \le 2c\,n\log_2^2 n$
con la misma constante, como exige el método.

\textbf{Cota inferior.} Conjetura: $T(n) \ge \frac{c}{2}\, n \log_2^2 n$ para todo $n \ge 1$.

\emph{Caso base, $n = 1$.} $T(1) = c$, y la cota vale
$\frac{c}{2} \cdot 1 \cdot 0 = 0$. Se cumple $c \ge 0$.

\emph{Paso inductivo.} Sea $n \ge 2$ y supongamos que la cota vale para $n/2$. Sustituyendo:
\[
\begin{aligned}
  T(n) &= 2\,T\!\left(\frac{n}{2}\right) + c\,n\log_2 n \\
       &\ge 2 \cdot \frac{c}{2}\, \frac{n}{2} \left(\log_2 n - 1\right)^2 + c\,n\log_2 n \\
       &= \frac{c}{2}\, n \left(\log_2^2 n - 2\log_2 n + 1\right) + c\,n\log_2 n \\
       &= \frac{c}{2}\, n\log_2^2 n + \frac{c}{2}\, n.
\end{aligned}
\]
Como $\frac{c}{2}\, n \ge 0$, queda $T(n) \ge \frac{c}{2}\, n\log_2^2 n$ con la misma
constante.

Juntando ambas cotas, para todo $n \ge 2$:
\[
  \frac{c}{2} \cdot n\log_2^2 n \le T(n) \le 2c \cdot n\log_2^2 n.
\]
Según la definición vista en clases, $T(n) \in \T{n\log^2 n}$ si existen dos constantes
$c_1, c_2 > 0$ y un $n_0 \in \N$ tales que $0 \le c_1 n\log_2^2 n \le T(n) \le c_2 n\log_2^2 n$
para todo $n \ge n_0$. Luego, $c_1 = c/2$ y $c_2 = 2c$, y las dos cotas valen desde
$n = 2$, así que $n_0 = 2$. Por lo tanto, $T(n) \in \T{n \log^2 n}$.

% ============================================================
\subsubsection*{(d) Mediana en tiempo lineal}

Si la mediana se pudiera obtener y los puntos se pudieran particionar en $\T{n}$, el
único cambio estaría en el trabajo no recursivo: el ordenamiento, de costo $\T{n\log n}$,
se reemplaza por una operación de costo $\T{n}$. Las dos llamadas recursivas siguen igual.
La nueva recurrencia es
\[
  T(n) = 2\,T\!\left(\frac{n}{2}\right) + c\,n.
\]
Es la misma recurrencia de Merge Sort, que se resolvió en clases: cada nivel del árbol
cuesta $c\,n$ y hay $\log_2 n + 1$ niveles, así que
\[
  T(n) \in \T{n\log n}.
\]
El costo de construcción baja de $\T{n\log^2 n}$ a $\T{n\log n}$, es decir, se ahorra un
factor $\log n$. Ese factor venía de ordenar en cada llamada, cuando para elegir la
mediana no hace falta ordenar todo. Una forma de lograrlo es el algoritmo de selección
con mediana de medianas, cuya recurrencia $T(n/5) + T(7n/10) + n$ se vio en clases y da
$\T{n}$.

